% Figure: streamlines, pathline, streakline in an unsteady flow, showing they
% differ when the flow is time-dependent.
\begin{figure}[htbp]
\centering
\begin{tikzpicture}[
  axis/.style={draw=engblack, -{Stealth}, thick},
  sline/.style={draw=enggray, thick},
  path/.style={draw=engblue, very thick},
  streak/.style={draw=engred, very thick, dashed},
  label/.style={font=\scriptsize},
]
\draw[axis] (-0.4,0) -- (9,0) node[anchor=west, font=\small] {$x$};
\draw[axis] (0,-2.4) -- (0,2.4) node[anchor=south, font=\small] {$y$};

% Instantaneous streamlines (snapshot, nearly horizontal at this instant)
\foreach \y in {-2,-1,1,2} {
  \draw[sline] (0.3,\y) -- (8.5,\y);
}
\node[label, enggray, anchor=west] at (8.6,2) {streamlines (snapshot)};

% Pathline of one particle in oscillating cross-flow:
% x grows steadily, y oscillates. (single particle from origin)
\draw[path] plot[domain=0:8, samples=120]
  ({\x}, {0.9*sin(\x*60)});
\node[label, engblue, anchor=south west] at (5.0,0.95) {pathline (one particle)};
\filldraw[engblue] (0,0) circle (1.3pt);

% Streakline: locus of particles released from origin, at one instant.
% Phase-shifted relative to the pathline.
\draw[streak] plot[domain=0:8, samples=120]
  ({\x}, {0.9*sin(\x*60 + 120)});
\node[label, engred, anchor=north west] at (5.0,-1.0) {streakline (dye trace)};

% Injection point
\node[label, anchor=north east] at (0,-0.1) {release point};

\end{tikzpicture}
\caption[Streamlines, pathline, and streakline in unsteady
flow]{The three flow-visualisation curves for a time-dependent flow
$u_x = U_0$, $u_y = V_0\cos(\omega t)$ with steady streamwise drift
and an oscillating cross-flow. The instantaneous streamlines (grey)
are tangent to the velocity field at one frozen instant. The pathline
(blue) is the trajectory of a single particle released at the origin.
The streakline (red dashed) is the locus, at one instant, of all
particles that have passed through the release point; it is the curve
a dye filament traces. The three curves differ because the flow is
unsteady. For a steady flow they collapse onto one another, and that
coincidence is what makes steady-flow visualisation easy to read.}
\label{fig:vol03ch11:streamline-types}
\end{figure}
